% Laget automatisk av elevark/lag_elevark.py fra opplegg/folg-grafen-sensor.html. Endringer her blir overskrevet.
\documentclass[11pt]{article}
\usepackage{elevark}
\begin{document}
\elevarkhode{Fysikk 1 · Rettlinjet bevegelse}{Følg grafen med bevegelsessensor}{https://fysikklab.no/opplegg/folg-grafen-sensor.html}{Opplegget på nett}

En bevegelsessensor måler hvor langt unna du står. Klarer du å gå slik at grafen for posisjon mot tid blir lik en gitt målgraf? Her må du oversette fra graf til bevegelse med hele kroppen: Hva betyr en vannrett linje, en bratt linje eller en krum kurve?

\section*{Hensikt}

\begin{itemize}
\item Koble formen på en x(t)-graf til hvordan man faktisk beveger seg.
\item Kjenne forskjellen på konstant fart og konstant akselerasjon i kroppen.
\end{itemize}

\section*{Utstyr}

\begin{itemize}
\item En bevegelsessensor (ultralyd) koblet til en logger eller datamaskin, med programvare som viser posisjonen mot tiden mens dere går. Sett målingen til 10~s.
\item Omtrent 3~m fri gulvplass rett foran sensoren, og en stor bok eller et pappstykke som den som går, holder foran magen. Da får sensoren en jevn flate å måle mot.
\item Uten sensor: Film den som går fra siden med en målestokk på gulvet, og analyser videoen i Tracker. Da tar det lengre tid.
\end{itemize}

\section*{Framgangsmåte}

\begin{enumerate}
\item Velg en målgraf i verktøyet på nettsiden. Skriv først ned med ord hvordan dere skal gå: hvor dere starter, hvilken vei, hvor fort og når dere stopper.
\item Still dere opp på startposisjonen. Én starter målingen, og den som går, prøver å følge grafen.
\item Kopier tabellen med tid og posisjon fra programvaren (eller Tracker), og lim den inn i verktøyet. Se hvor dere bommet.
\item Prøv igjen og forsøk å få avviket ned. Bytt så rolle og målgraf.
\end{enumerate}

\section*{Oppgaver}

\begin{oppgaveliste}
\oppgave \fase{gjett}{Gjett}Se på målgraf $B$ før dere prøver. Beskriv med ord hvordan dere må gå. Hvor fort må dere gå i hver del?
\svar
\oppgave \fase{test}{Test}Prøv målgraf $A$, $B$ og $E$. Hvilken del av grafene var vanskeligst å treffe?
\svar
\oppgave \fase{forklar}{Forklar}Hva betyr det når grafen er vannrett? Når den er bratt? Hva betyr det at grafen går nedover?
\svar
\oppgave \fase{test}{Test}Prøv målgraf $C$ og $D$. Hvordan må farten endre seg underveis? Hva er forskjellen på dem?
\svar
\oppgave \fase{test}{Tegn}Tegn v(t)-grafen som hører til målgraf $E$. Hvor stor er farten i hver del?
\svar
\oppgave \fase{utfordring}{Utfordring}Kan noen gå slik at x(t)-grafen får en loddrett del? Lag en egen målgraf på papir som en annen gruppe skal prøve å følge.
\svar
\end{oppgaveliste}

\end{document}
